\documentclass[11pt,a4paper]{article}
\usepackage[T1]{fontenc}
\usepackage[utf8]{inputenc}
\usepackage{lmodern}
\usepackage{microtype}
\usepackage[margin=27mm,headheight=15pt]{geometry}
\usepackage{amsmath,amssymb,amsthm,mathtools}
\usepackage{booktabs,array,longtable,enumitem}
\usepackage{xcolor}
\usepackage{fancyhdr}
\usepackage{hyperref}
\usepackage{bookmark}
\usepackage{xurl}
\definecolor{inkblue}{HTML}{183D59}
\definecolor{linkblue}{HTML}{215F87}
\hypersetup{colorlinks=true,linkcolor=linkblue,citecolor=linkblue,urlcolor=linkblue,
 pdftitle={All-Degree Rigidity of a Weighted Keller Class},
 pdfauthor={ChatGPT; prepared for Vladimir Reshetnikov},
 pdfsubject={Classification, degenerations, and an obstructed double point}}
\pagestyle{fancy}
\fancyhf{}
\fancyhead[L]{\small\color{inkblue}All-degree rigidity of a weighted Keller class}
\fancyhead[R]{\small Research draft}
\fancyfoot[C]{\thepage}
\renewcommand{\headrulewidth}{0.3pt}
\setlength{\parskip}{0.35em}
\setlength{\parindent}{1.2em}
\setlength{\emergencystretch}{2em}
\setlist{itemsep=0.25em,topsep=0.4em}
\numberwithin{equation}{section}
\newtheorem{theorem}{Theorem}[section]
\newtheorem{lemma}[theorem]{Lemma}
\newtheorem{proposition}[theorem]{Proposition}
\newtheorem{corollary}[theorem]{Corollary}
\theoremstyle{definition}
\newtheorem{definition}[theorem]{Definition}
\newtheorem{example}[theorem]{Example}
\theoremstyle{remark}
\newtheorem{remark}[theorem]{Remark}
\newcommand{\Q}{\mathbb{Q}}
\newcommand{\A}{\mathbb{A}}
\newcommand{\Gm}{\mathbb{G}_{\mathrm m}}
\newcommand{\J}{\mathcal{J}}
\newcommand{\Spec}{\operatorname{Spec}}
\newcommand{\diag}{\operatorname{diag}}
\newcommand{\wt}{\operatorname{wt}}
\newcommand{\coeff}[2]{[#1]#2}
\newcommand{\cstar}{c_*}
\newcommand{\eps}{\varepsilon}
\newcommand{\repo}{\texttt{ProveIt}}
\newcommand{\dd}{\mathrm d}

\begin{document}
\begin{titlepage}
\thispagestyle{empty}
\noindent{\color{inkblue}\large RESEARCH REPORT}\par
\vspace{1.7cm}
{\color{inkblue}\Huge\bfseries All-Degree Rigidity\par
\vspace{0.2cm}of a Weighted Keller Class\par}
\vspace{0.65cm}
{\Large Classification, degenerations,\par and an obstructed double point\par}
\vspace{1.1cm}
{\large Prepared for Vladimir Reshetnikov}\par
\medskip
September 24, 2026\par
\vspace{0.9cm}
\begin{abstract}
We classify, without any bound on ordinary degree, a class of polynomial
Keller maps in three variables naturally singled out by the Jacobian project
in the \repo\ repository. The source weights are $(-1,1,2)$ and the target
weights are $(2,1,-1)$. In invariant coordinates $t=xy$ and $v=x^2z$, we
assume that the first two numerators are affine in $v$ and that the third
numerator is affine in $(t,v)$. Every normalized map is either an explicitly
invertible tame map or a member of a two-parameter diagonal orbit of the
known Alp\"oge map. In the latter case its degree profile is forced to be
$(7,6,4)$ and its expanded support has exactly sixteen monomials.

The proof replaces a degree-bounded coefficient search by one-variable
polynomial rigidity. An equivariant source shear extends the classification
 to arbitrary polynomial dependence of $r$ on $t$ when its $v$-coefficient
 is constant, and gives the exact noninvertible degree spectrum
 $\{7,8,10,12,\ldots\}$. We also identify the reduced closure of the nonconstant
branch as an affine plane whose coordinate axes are precisely its
invertible members. Finally, for the normalized degree-seven coefficient
problem with third numerator $1+t+v$, we prove that the entire coefficient
scheme is a double point, $\Spec\Q[\eps]/(\eps^2)$. An explicit nonzero
infinitesimal deformation is obstructed at second order. All finite
identities and obstruction certificates are independently reproduced by an
included exact-arithmetic script.
\end{abstract}
\vfill
\noindent\small
\textbf{Status.} Human-readable proofs and reproducible symbolic certificates;
not a Lean/Rocq formalization or a peer-reviewed publication. The results
extend the inspected repository statements, but publication priority is not
established. This report does not claim a new counterexample to the Jacobian
conjecture or a classification of unrestricted Keller maps.
\end{titlepage}

\begingroup
\footnotesize
\setlength{\parskip}{0pt}
\tableofcontents
\endgroup
\newpage

\section{Research question, provenance, and contribution}
\label{sec:intro}

A useful research question should distinguish a structural theorem from a
large but finite search. The Jacobian component of \repo\ provides a
particularly concrete opportunity: it gives explicit maps, exact collision
certificates, a weighted determinant reduction, and degree-bounded searches
for simpler representatives. We ask whether one of its natural coefficient
classes is rigid \emph{without imposing a degree bound}. We then ask what
remains of the apparent uniqueness when coefficients are allowed to contain
nilpotents. These are separate questions: a unique field-valued solution
need not be a reduced point of its coefficient scheme.

\subsection{The repository baseline}

The inspected snapshot is commit
\begin{center}
\texttt{e21766d04c2b8a9b2cdba0cd43e563024b3bd1b9}.
\end{center}
The relevant files are the Jacobian project README, its research README, and
\texttt{verify\_equivariant\_sparsity.py}; immutable references are supplied
in the bibliography \cite{repo-main,repo-research,repo-sparsity}.

The repository's three-variable map can be written
\begin{equation}\label{eq:original}
\begin{gathered}
u=1+xy,\qquad h=u^2z+y^2(1+3u),\\
F_0(x,y,z)=\bigl(uh,\ y+3xh,\ x(5-3u-x^2z)\bigr).
\end{gathered}
\end{equation}
Its Jacobian determinant is $-2$, and
\begin{equation}\label{eq:oldcollision}
F_0(-1,1,5)=F_0(0,-2,-16)=(0,-2,0).
\end{equation}
These identities, their formal verification, and their provenance are
already part of the repository \cite{repo-main}. We verify the displayed
identities again in the companion script but do not present them as new.

The research files restrict a sparsity calculation to maps of ordinary
degree at most seven whose invariant numerators $p,q$ are affine in $v$ and
whose third numerator $r$ is affine in $(t,v)$. After normalizing
$r=1+t+v$, the script proves that eleven optional coefficients cannot vanish,
and identifies the branch $q_{40}=0$ with the normalized map. Its certificate
contains a square $(q_{21}-4)^2$, but the repository discussion does not
identify the entire coefficient scheme as the dual numbers
\cite{repo-research,repo-sparsity}.

\subsection{What this article establishes}

The main result, Theorem~\ref{thm:classification}, removes the ordinary
degree restriction from that structural class. It also handles constant
$r$ and proves that the two one-term nonconstant cases are impossible. The
proof forces the normalized nonconstant map outright, so it proves
$q_{40}=0$ rather than imposing it.

Theorem~\ref{thm:triangular} extends the result beyond affine $r$ by an
explicit triangular normal form. Its degree-spectrum corollary determines
all possible maximum degrees in that larger class.
Theorem~\ref{thm:plane} describes the reduced two-parameter closure and its
invertible boundary. Theorem~\ref{thm:double} strengthens uniqueness in the
fixed normalized degree-seven slice to an exact scheme presentation, and Theorem~\ref{thm:torusdouble} identifies the full
scheme on the nonconstant-parameter open locus. These results connect
three levels of information:
\begin{center}
\begin{tabular}{@{}p{0.25\textwidth}p{0.66\textwidth}@{}}
\toprule
Level & Conclusion proved here\\
\midrule
Field-valued maps & An all-degree automorphism/counterexample dichotomy.\\
Reduced families & No deformation after fixing $r=1+t+v$; a parameter plane
before fixing its two nonconstant coefficients.\\
Nonreduced families & One square-zero parameter in the degree-seven slice,
with a nonzero second-order obstruction.\\
\bottomrule
\end{tabular}
\end{center}

\subsection{Relation to contemporary literature and limits}

Shaska studies graded Keller maps, their quotient geometry, and the role of
the sign pattern of weights \cite{shaska}. Gao gives a tangent-sweep
construction and further explicit examples, including a three-dimensional
map of degree four \cite{gao}. Our subject is narrower: rigidity within a
specified invariant-numerator class, not the construction of the known map
or a minimum-degree theorem for all three-variable maps. In particular, our
forced degree seven is compatible with examples outside this class of lower
degree.

The searched sources and the inspected repository justify the comparison
above, not a universal priority claim. The mathematical contribution of this
report is its stated and proved classification and deformation results.
Whether equivalent formulations already occur elsewhere requires further
bibliographic review. Nothing in the proofs depends on accepting a new
external claim about the general Jacobian problem: the particular map,
its determinant, and its collisions are explicit.

\section{Weighted coordinates and the exact determinant reduction}
\label{sec:coordinates}

Throughout the classification, $K$ is an arbitrary field of characteristic
zero. All polynomial equalities are formal equalities over $K$.

\begin{definition}
A polynomial map $F:\A^3_K\to\A^3_K$ is a \emph{Keller map} if its formal
Jacobian determinant is a nonzero element of $K$. We use the source action
\[
(x,y,z)\longmapsto(\rho^{-1}x,\rho y,\rho^2z)
\]
and the target action with weights $(2,1,-1)$. Equivariance means that the
three coordinate polynomials are weighted homogeneous of weights $2,1,-1$,
respectively.
\end{definition}

\subsection{Invariant numerators}

Set
\begin{equation}\label{eq:invariants}
t=xy,\qquad v=x^2z.
\end{equation}
Both are invariant under the source action. Every equivariant polynomial
map has a unique expression in the localization at $x$ of the form
\begin{equation}\label{eq:lift}
F=\mathcal L(p,q,r)
 :=\left(\frac{p(t,v)}{x^2},\frac{q(t,v)}x,xr(t,v)\right),
\qquad p,q,r\in K[t,v].
\end{equation}
For example, a monomial $x^iy^jz^k$ of weight $2$ satisfies
$i=j+2k-2$; multiplying it by $x^2$ gives $t^jv^k$.
The corresponding formulas for weights $1$ and $-1$ are
$i=j+2k-1$ and $i=j+2k+1$. They prove existence and uniqueness of the
numerators monomial by monomial.

For general numerators, polynomiality requires that every monomial $t^iv^j$
in $p$ satisfy $i+2j\ge2$, and every such monomial in $q$ satisfy
$i+2j\ge1$. In the class studied here, $\deg_vp,\deg_vq\le1$, so these
conditions reduce exactly to
\begin{equation}\label{eq:polynomiality}
p(0,0)=p_t(0,0)=q(0,0)=0.
\end{equation}
There is no additional restriction on $r$ for polynomiality. The apparent
fractions in \eqref{eq:lift} are therefore polynomial expressions whenever
\eqref{eq:polynomiality} holds in our class.

\begin{proposition}[Weighted determinant identity]\label{prop:jacobian}
For a polynomial lift \eqref{eq:lift},
\begin{equation}\label{eq:J}
\det J_F=\J(p,q,r):=
\det\begin{pmatrix}
-2p&p_t&p_v\\
-q&q_t&q_v\\
r&r_t&r_v
\end{pmatrix},
\end{equation}
where the right side is evaluated at $(t,v)=(xy,x^2z)$.
\end{proposition}
\begin{proof}
Compute in coordinates $(x,t,v)$ on $x\ne0$. The coordinate change
$(x,y,z)\mapsto(x,xy,x^2z)$ has determinant $x^3$. The derivative matrix of
$(x^{-2}p,x^{-1}q,xr)$ with respect to $(x,t,v)$ has rows
\[
(-2x^{-3}p,x^{-2}p_t,x^{-2}p_v),\quad
(-x^{-2}q,x^{-1}q_t,x^{-1}q_v),\quad
(r,xr_t,xr_v).
\]
Factoring $x^{-2},x^{-1},x$ from its rows and $x^{-1}$ from its first column
leaves the matrix in \eqref{eq:J} and contributes $x^{-3}$. The chain rule
cancels this with $x^3$. Localization is injective on $K[x,y,z]$, so the
identity holds as a polynomial identity, including on $x=0$.
\end{proof}

\subsection{Normalization is exhaustive}

Suppose that \eqref{eq:polynomiality} holds and $\J$ is a nonzero constant.
Evaluating at $(t,v)=(0,0)$ gives
\begin{equation}\label{eq:originJ}
\J(0,0)=-p_v(0,0)q_t(0,0)r(0,0).
\end{equation}
All three factors are therefore nonzero. Dividing the three target
coordinates by these factors normalizes
\begin{equation}\label{eq:normalization}
p_v(0,0)=q_t(0,0)=r(0,0)=1.
\end{equation}
For a Keller map, the resulting constant determinant is $-1$.
This operation preserves the structural restrictions on $p,q,r$, ordinary
degrees, support sizes, and invertibility or noninvertibility. Thus a
classification in this normalization classifies the unnormalized maps up
to independent nonzero target scalings. We can write
\begin{equation}\label{eq:r}
r=1+\beta t+\gamma v.
\end{equation}
The word \emph{normalized} below always includes
\eqref{eq:polynomiality} and \eqref{eq:normalization}; it does not implicitly
assume $\beta=\gamma=1$.

\section{The complete all-degree classification}
\label{sec:main}

Define the following polynomials over $\Q[a,b,t,v]$:
\begin{subequations}\label{eq:family}
\begin{align}
p_{a,b}={}&v-\frac{8a}{9}t^2-2ab\,tv
 +\frac{28a^2b}{27}t^3+\frac{4a^2b^2}{3}t^2v \notag\\
 &\hspace{1.2em}-\frac{8a^3b^2}{27}t^4-\frac{8a^3b^3}{27}t^3v,
 \label{eq:familyp}\\
q_{a,b}={}&t+9bv-8ab\,t^2-12ab^2tv
 +4a^2b^2t^3+4a^2b^3t^2v,\label{eq:familyq}\\
r_{a,b}={}&1+ab\,t+ab^2v.\label{eq:familyr}
\end{align}
\end{subequations}
Let $F_{a,b}=\mathcal L(p_{a,b},q_{a,b},r_{a,b})$.

\begin{theorem}[All-degree classification]\label{thm:classification}
Let $p,q,r\in K[t,v]$ satisfy $\deg_vp,\deg_vq\le1$, let $r$ be affine in
$(t,v)$, and suppose that their lift is polynomial and normalized as in
\eqref{eq:normalization}. Then $\J(p,q,r)=-1$ if and only if exactly one
of the following descriptions applies.
\begin{enumerate}[label=\textup{(\Alph*)}]
\item\label{case:tame} $r=1$ and, for unique $\lambda\in K$ and
$B\in t^2K[t]$,
\begin{equation}\label{eq:tamepq}
p=v+B(t),\qquad q=t+\lambda\bigl(v+B(t)\bigr).
\end{equation}
Its lift is a tame polynomial automorphism.
\item\label{case:nonconstant} $r$ is nonconstant, both $\beta$ and $\gamma$
in \eqref{eq:r} are nonzero, and
\begin{equation}\label{eq:abrecover}
(p,q,r)=(p_{a,b},q_{a,b},r_{a,b}),\qquad
 a=\frac{\beta^2}{\gamma},\quad b=\frac{\gamma}{\beta}.
\end{equation}
Its lift is noninjective. Its coordinate degrees are $(7,6,4)$, and its
coordinate support sizes are $(7,6,3)$.
\end{enumerate}
In particular, there are no normalized Keller maps in this class with
exactly one of $\beta,\gamma$ nonzero. No ordinary degree bound is assumed.
\end{theorem}

We prove the classification in Sections~\ref{sec:lemmas}--\ref{sec:proof}.
The invertibility, collisions, and support assertions are proved in
Section~\ref{sec:consequences}. The forward direction is the substantive
rigidity statement; the reverse direction follows either from the formulas
there or by substitution into \eqref{eq:J}.

\subsection{The normalized representative}

Put $F_*=F_{1,1}$. Its numerators are
\begin{align*}
p_*&=v-\frac89t^2-2tv+\frac{28}{27}t^3
       +\frac43t^2v-\frac8{27}t^4-\frac8{27}t^3v,\\
q_*&=t+9v-8t^2-12tv+4t^3+4t^2v,\qquad r_*=1+t+v.
\end{align*}
It is exactly a normalization of \eqref{eq:original}:
\begin{equation}\label{eq:originalnormalization}
F_*=\diag\left(-\frac12,-\frac32,\frac12\right)
\circ F_0\circ\diag\left(1,-\frac23,-2\right).
\end{equation}
Thus the nonconstant part of the classification identifies a known
representative rather than producing an unrelated counterexample.

\section{One-variable rigidity lemmas}
\label{sec:lemmas}

The decisive step is that constant weighted Wronskians control degree even
when arbitrarily high-degree polynomials are allowed initially.

\begin{lemma}[A degree bound from a constant weighted Wronskian]
\label{lem:wronskian}
Let $B,D\in K[t]$ satisfy
\begin{equation}\label{eq:wronskian}
DB'-2BD'=\kappa,\qquad \kappa\in K^\times.
\end{equation}
Then $D\ne0$, $\deg D\le1$, and $\deg B\le2$. More precisely, either $D$
is a nonzero constant and $B$ has degree one, or $D$ has degree one and
$B=\eta D^2+\mu$ for some $\eta,\mu\in K$, with
$-2\mu D'=\kappa$.
\end{lemma}
\begin{proof}
Neither $B$ nor $D$ can be zero. If $D$ is constant, \eqref{eq:wronskian}
gives $B'=\kappa/D$, so $B$ has degree one.

Now put $n=\deg D\ge1$ and $m=\deg B$. Unless $m=2n$, the leading term
on the left has degree $m+n-1$ and nonzero coefficient proportional to
$m-2n$. This uses characteristic zero. If $m>2n$, it contradicts the
constant right-hand side. If $m=2n$, subtract from $B$ the unique scalar
multiple of $D^2$ that cancels its leading coefficient. This changes
neither side of \eqref{eq:wronskian}. The resulting polynomial $B_0$ is
nonzero and has degree $m_0<2n$. If initially $m<2n$, just take $B_0=B$.

The leading term calculation for $B_0$ now gives $m_0+n-1=0$. As
$m_0\ge0$ and $n\ge1$, we obtain $n=1$ and $m_0=0$. Hence
$B=\eta D^2+\mu$, and substitution gives $-2\mu D'=\kappa$.
\end{proof}

\begin{lemma}[A cubic--square relation]\label{lem:powers}
Suppose $A,C\in K[t]$ are nonzero and
\begin{equation}\label{eq:powers}
2CA'-3AC'=0,\qquad A(0)=1.
\end{equation}
Then there are unique $E\in K[t]$ with $E(0)=1$ and $c\in K^\times$ such
that
\begin{equation}\label{eq:EC}
A=E^3,\qquad C=cE^2.
\end{equation}
\end{lemma}
\begin{proof}
The derivative of $A^2/C^3$ is zero. A rational function over a
characteristic-zero field with derivative zero lies in $K$: writing it in
coprime numerator--denominator form, the identity $f'g-fg'=0$ forces each
nonconstant numerator or denominator to divide its own derivative, which
is impossible by degree. Thus $A^2/C^3\in K^\times$.

For every irreducible polynomial $\pi$, unique factorization gives
$2\operatorname{ord}_\pi A=3\operatorname{ord}_\pi C$. Consequently the
multiplicities in $A$ are multiples of three and those in $C$ are the
corresponding multiples of two. Because $A(0)=1$, none of the common
factors vanishes at zero. Normalize their product to have value one there.
The resulting polynomial $E$ satisfies $A=E^3$ and $C=C(0)E^2$.
No extraction of a scalar root or extension of $K$ is necessary.
Uniqueness follows from the factorization and the value at zero.
\end{proof}

\begin{lemma}[The remaining common factor is linear]\label{lem:Edegree}
Suppose that $A(0)=1$, $B(0)=-1$, and
\begin{align}
DB'-2BD'&=-1,\label{eq:E0}\\
DA'+2CB'-3AD'-2BC'&=0,\label{eq:E1}\\
2CA'-3AC'&=0.\label{eq:E2}
\end{align}
If $C\ne0$, then in \eqref{eq:EC} the polynomial $E$ has degree at most
one.
\end{lemma}
\begin{proof}
Lemma~\ref{lem:wronskian} bounds $\deg D\le1$ and $\deg B\le2$.
Lemma~\ref{lem:powers} gives $A=E^3$, $C=cE^2$. Suppose $m=\deg E\ge2$.

If $D$ has degree one, the leading coefficient of $DA'-3AD'$ in degree
$3m$ is $3(m-1)$ times the leading coefficient of $D$ times the cube of
that of $E$. It is nonzero. The other part $2CB'-2BC'$ of \eqref{eq:E1}
has degree at most $2m+1<3m$, a contradiction.

If $D$ is constant, it is nonzero and $B$ has degree one. The term $DA'$
has degree $3m-1$, whereas $2CB'-2BC'$ has degree at most $2m$.
Again the leading term cannot cancel. Thus $\deg E\le1$.
\end{proof}

\begin{remark}
The lemmas address the possible high-degree cancellations explicitly. A
coefficient search stopped at any fixed degree would not establish these
bounds. Characteristic zero is essential to the leading coefficient
arguments and to the assertion about rational functions with zero derivative.
\end{remark}

\section{Proof of the classification, including exceptional cases}
\label{sec:proof}

\subsection{\texorpdfstring{The case $\gamma\ne0$: a moving coordinate}{The case gamma nonzero: a moving coordinate}}

First reduce $r=1+\beta t+\gamma v$ to coefficient one on $v$. Replace
\begin{equation}\label{eq:gammachange}
\widetilde p(t,v)=\gamma p(t,v/\gamma),\quad
\widetilde q(t,v)=q(t,v/\gamma),\quad
\widetilde r(t,v)=r(t,v/\gamma).
\end{equation}
The determinant and normalization are unchanged, while
$\widetilde r=1+\beta t+v$. We temporarily omit the tildes.

Set $s=1+\beta t+v$ and write uniquely
\begin{equation}\label{eq:ABCD}
p=A(t)s+B(t),\qquad q=C(t)s+D(t),\qquad r=s.
\end{equation}
Substitution into \eqref{eq:J} gives
\begin{equation}\label{eq:threeidentities}
\J=E_0+E_1s+E_2s^2,
\end{equation}
where the expressions $E_0,E_1,E_2$ are the left sides of
\eqref{eq:E0}--\eqref{eq:E2}. For clarity, differentiation at fixed $v$
gives $p_t=A's+A\beta+B'$, $q_t=C's+C\beta+D'$, and $r_t=\beta$.
Expanding the determinant cancels the $\beta$ terms and yields
\eqref{eq:threeidentities}. Since $(t,s)$ are independent polynomial
coordinates, the Keller equation is exactly
\eqref{eq:E0}--\eqref{eq:E2}.

The boundary conditions become
\begin{equation}\label{eq:boundaryABCD}
\begin{gathered}
A(0)=1,\quad B(0)=-1,\quad A'(0)+\beta+B'(0)=0,\\
C(0)+D(0)=0,\qquad C'(0)+\beta C(0)+D'(0)=1.
\end{gathered}
\end{equation}
We must not divide by $C$ before excluding $C=0$. If $C$ were zero, the
fourth condition would give $D(0)=0$. Equation~\eqref{eq:E0} at zero would
then give $D'(0)=-1/2$, while \eqref{eq:E1} at zero would give
$-3D'(0)=0$. This is impossible.

The rigidity lemmas now apply. Write
\begin{equation}\label{eq:linearforms}
\begin{gathered}
E=1+kt,\quad A=E^3,\quad C=cE^2,\qquad c\ne0,\\
D=dt-c,\qquad B=-1+b_1t+b_2t^2.
\end{gathered}
\end{equation}
The first derivative condition in \eqref{eq:boundaryABCD} gives
\begin{equation}\label{eq:b1}
b_1=-\beta-3k.
\end{equation}
After dividing \eqref{eq:E1} by the nonzero polynomial $E$, it reads
\begin{equation}\label{eq:keylinear}
2c(EB'-2kB)=3(kc+d)E.
\end{equation}
Its constant and linear coefficients, together with \eqref{eq:b1}, give
\begin{equation}\label{eq:bd}
b_2=-\beta k-2k^2,\qquad
\frac dc=-\frac{2\beta+5k}{3}.
\end{equation}
The final condition in \eqref{eq:boundaryABCD} is
$c(\beta+2k)+d=1$, which therefore simplifies to
\begin{equation}\label{eq:ck}
c(\beta+k)=3.
\end{equation}
The coefficient of $t$ in \eqref{eq:E0} is $-db_1-2cb_2$. Its vanishing,
using \eqref{eq:b1} and \eqref{eq:bd}, is
\[
\frac c3(\beta+k)(2\beta+3k)=0.
\]
Equation~\eqref{eq:ck} permits cancellation of the first factors. Hence
\begin{equation}\label{eq:endpoint}
k=-\frac{2\beta}{3},\qquad
\beta\ne0,\qquad c=\frac9\beta,\qquad d=4,
\qquad b_1=\beta,\qquad b_2=-\frac{2\beta^2}{9}.
\end{equation}
The remaining constant coefficient of \eqref{eq:E0} is indeed $-1$:
$-cb_1+2d=-9+8=-1$. All three equations hold for these values. We have proved
uniqueness and existence, including the exclusion of $\beta=0$.

More compactly, before undoing \eqref{eq:gammachange}, the answer is
\begin{equation}\label{eq:compactendpoint}
\begin{gathered}
E=1-\frac{2\beta}{3}t,\qquad s=1+\beta t+v,\\
p=E^3s-1+\beta t-\frac{2\beta^2}{9}t^2,\qquad
q=\frac9\beta E^2s+4t-\frac9\beta.
\end{gathered}
\end{equation}
Replace $v$ by $\gamma v$ and divide $p$ by $\gamma$. Expansion yields
\eqref{eq:family} with $a=\beta^2/\gamma$ and $b=\gamma/\beta$.
This proves the nonconstant branch whenever $\gamma\ne0$.

\subsection{\texorpdfstring{The case $\gamma=0$, $\beta\ne0$: no solutions}{The case gamma zero, beta nonzero: no solutions}}

Write $p=A(t)v+B(t)$, $q=C(t)v+D(t)$, and $r=1+\beta t$. The normalization
implies $A(0)=1$ and $B(0)=B'(0)=D(0)=0$. The coefficient of $v$ in
\eqref{eq:J} is
\begin{equation}\label{eq:gamma0v}
\beta AC+r(A'C-AC')=0.
\end{equation}
If $C\ne0$, this says that $(rA/C)'=0$, so $C=\lambda rA$ for some
$\lambda\in K^\times$. When $C=0$, the same expression holds with
$\lambda=0$. The constant coefficient of the determinant then becomes
\begin{equation}\label{eq:gamma0const}
A\left(\lambda(r^2B)'-(rD)'\right)=-1.
\end{equation}
A product of two polynomials is a nonzero constant only when both are
constants. Thus $A=1$. Integrating \eqref{eq:gamma0const} and using
$B(0)=D(0)=0$ gives
\begin{equation}\label{eq:impossible}
\lambda r^2B-rD=-t.
\end{equation}
Evaluate at $t=-1/\beta$. The left side is zero and the right side is
$1/\beta\ne0$. This contradiction excludes the entire case, at every degree.

\subsection{\texorpdfstring{The case $r=1$: the full tame branch}{The case r=1: the full tame branch}}

Again write $p=Av+B$ and $q=Cv+D$. The coefficient of $v$ in the determinant
is $A'C-AC'=0$. Since $A(0)=1$, it follows that $C=\lambda A$ for some
$\lambda\in K$, allowing $C=0$. The constant coefficient is
\[
A(\lambda B'-D')=-1.
\]
Consequently $A=1$ and $D=t+\lambda B$, using $B(0)=D(0)=0$.
Polynomiality and normalization give $B(0)=B'(0)=0$, hence $B\in t^2K[t]$.
This is precisely \eqref{eq:tamepq}. Conversely, these formulas give
$\J=-1$ directly. Together with the preceding two cases, this completes
the algebraic classification in Theorem~\ref{thm:classification}.

\section{Invertibility, diagonal equivalence, and sharp support}
\label{sec:consequences}

\subsection{An inverse for every constant-third-numerator map}

Write $B(t)=t^2H(t)$. The lift of \eqref{eq:tamepq} is
\begin{equation}\label{eq:tameF}
F(x,y,z)=\bigl(P,\ y+\lambda xP,\ x\bigr),\qquad
P=z+y^2H(xy).
\end{equation}
For output coordinates $(U,V,W)$, define
\begin{equation}\label{eq:tameinverse}
 x=W,\qquad y=V-\lambda WU,\qquad
 z=U-(V-\lambda WU)^2H\bigl(W(V-\lambda WU)\bigr).
\end{equation}
Substitution in both directions is immediate from $P=U$. This proves
polynomial invertibility. More precisely, \eqref{eq:tameF} is the
composition of the elementary shear
$z\mapsto z+y^2H(xy)$, the elementary shear $y\mapsto y+\lambda xz$ in the
updated coordinates, and the coordinate permutation
$(x,y,z)\mapsto(z,y,x)$. It is therefore tame. There is no bound on
$\deg H$: arbitrarily large degrees occur on the invertible branch.

\subsection{Every nonconstant map is in the same diagonal orbit}

For $a,b\ne0$, inspection of \eqref{eq:family} gives
\begin{equation}\label{eq:diagonal}
F_{a,b}=\diag\bigl((ab^2)^{-1},(ab)^{-1},1\bigr)
\circ F_*\circ\diag(1,ab,ab^2).
\end{equation}
The product of the two diagonal determinants is one, consistently with
$\det J_{F_{a,b}}=-1$. Thus there are two parameters before taking this
source--target equivalence, and no parameter after taking it. The scalar
parameters are uniquely recoverable from either $r$ or two coefficients:
\begin{equation}\label{eq:coeffrecover}
a=-\frac98\coeff{t^2}{p},\qquad
b=\frac19\coeff{v}{q}.
\end{equation}

\begin{corollary}[Sharp all-degree support rigidity]\label{cor:support}
Within the class of Theorem~\ref{thm:classification}, every noninvertible
map has ordinary maximum degree seven and exactly sixteen expanded
monomials. After ordering coordinates by their prescribed weights, its
degree profile is $(7,6,4)$ and its support profile is $(7,6,3)$.
\end{corollary}
\begin{proof}
Only branch~\ref{case:nonconstant} can be noninvertible. For $a,b\ne0$,
all coefficients displayed in \eqref{eq:family} are nonzero. Different
monomials in a numerator lift to different monomials. The degree of a lifted
$t^iv^j$ is $2i+3j-2$ in the first coordinate, $2i+3j-1$ in the second,
and $2i+3j+1$ in the third. The asserted profiles follow. Noninjectivity of
each map is established just below.
\end{proof}

This is stronger than saying that sixteen terms are minimal under a degree
cap: every noninvertible map in the structural class has that precise
support profile, regardless of the degree initially allowed. It is not a
statement about minimum degree or support outside the class.

\subsection{Collisions valid over the ground field}

Let $\gamma=ab^2$ and $\ell=9b$. Direct evaluation on the plane $x=0$
gives
\begin{equation}\label{eq:xzero}
F_{a,b}(0,y,z)=\left(z-\frac{8a}{9}y^2,y,0\right).
\end{equation}
At $t=0$, the numerators simplify to $p=v$, $q=\ell v$, and
$r=1+\gamma v$. Therefore
\begin{equation}\label{eq:newcollision}
F_{a,b}\left(1,0,-\frac1\gamma\right)
=F_{a,b}\left(0,-\frac\ell\gamma,\frac{71}\gamma\right)
=\left(-\frac1\gamma,-\frac\ell\gamma,0\right).
\end{equation}
To check the constant $71$, note that
$-(8a/9)\ell^2=-72\gamma$; the first coordinate of the second point's
image is therefore $(71-72)/\gamma$.

The two source points have different $x$ coordinates. This proves
noninjectivity over \emph{every} characteristic-zero ground field $K$,
without passing to an algebraic closure. It also gives a rational
certificate when $a,b\in\Q^\times$.

Equivariance propagates this to every $\rho\in K^\times$:
\begin{equation}\label{eq:collisionfamily}
\begin{aligned}
F_{a,b}\left(\rho,0,-\frac1{\gamma\rho^2}\right)
&=F_{a,b}\left(0,-\frac\ell{\gamma\rho},
                         \frac{71}{\gamma\rho^2}\right)\\
&=\left(-\frac1{\gamma\rho^2},-\frac\ell{\gamma\rho},0\right).
\end{aligned}
\end{equation}
These are explicit certificates in the coordinates of this report. They
are consequences of the classified family, not a claim that the collision
mechanism of the original map was unknown.

\begin{example}[The normalized rational witness]
For $a=b=1$, the points $(1,0,-1)$ and $(0,-9,71)$ have common image
$(-1,-9,0)$. The Jacobian is nonsingular at both points because it has
constant determinant $-1$. Thus local nonsingularity and global injectivity
are genuinely different conditions in the displayed family.
\end{example}

\subsection{A larger class: arbitrary dependence on the first invariant}
\label{sec:triangular}

The affine restriction on $r$ can be relaxed substantially without changing
its essential normal form. The relevant source automorphisms are
\begin{equation}\label{eq:source-shear}
T_h(x,y,z)=\bigl(x,y,z+y^2h(xy)\bigr),\qquad h\in K[t].
\end{equation}
They preserve the weights, have Jacobian one, and induce
$(t,v)\mapsto(t,v+t^2h(t))$. Their inverses are $T_{-h}$.

\begin{theorem}[All-degree triangular extension]\label{thm:triangular}
Let $p,q$ be affine in $v$ and let
\[
r=R(t)+\gamma v,\qquad R\in K[t],\quad \gamma\in K,
\]
with polynomial lift and normalized origin coefficients. In particular,
$R(0)=1$. Then the normalized lift is Keller if and only if one of the
following holds.
\begin{enumerate}[label=\textup{(\roman*)}]
\item $\gamma=0$, $R=1$, and the map is a tame map
\eqref{eq:tameF}.
\item $\gamma\ne0$, $\beta:=R'(0)\ne0$, and the map is
\begin{equation}\label{eq:triangular-normal}
F=F_{a,b}\circ T_h,\qquad
 a=\frac{\beta^2}{\gamma},\quad b=\frac{\gamma}{\beta},\quad
 h(t)=\frac{R(t)-1-\beta t}{\gamma t^2}.
\end{equation}
Here $h$ is a polynomial, uniquely determined by $r$. The map is
noninjective over $K$.
\end{enumerate}
Thus every noninvertible map in this larger class is the classified
representative up to diagonal scalings and one explicit equivariant source
shear. There is no degree bound on $R$, $p$, or $q$.
\end{theorem}
\begin{proof}
When $\gamma\ne0$, the numerator defining $h$ vanishes to order at least
two at zero. Composing the source with $T_{-h}$ replaces $v$ by
$v-t^2h(t)$ and changes the third numerator to $1+\beta t+\gamma v$.
The first two numerators remain affine in $v$; polynomiality and all
normalized origin coefficients are preserved. The source shear has
Jacobian one. Theorem~\ref{thm:classification} therefore applies: it
forces $\beta\ne0$ and identifies $F\circ T_{-h}$ with $F_{a,b}$.
This is \eqref{eq:triangular-normal}. The converse follows from the same
composition, and collision points are transported by $T_{-h}$.

When $\gamma=0$, use the proof of Section~\ref{sec:proof} with
$r=R(t)$ instead of $1+\beta t$. The coefficient of $v$ is
$r'AC+r(A'C-AC')=0$, whence $C=\lambda rA$, including $C=0$.
The constant equation is still
$A\bigl(\lambda(r^2B)'-(rD)'\bigr)=-1$.
It forces $A=1$ and, on integration with the origin conditions,
$\lambda r^2B-rD=-t$. Thus $r$ divides $t$. Since $r(0)=1$, this
forces $r$ to be constant and hence $r=1$. The already proved tame
classification finishes this case.
\end{proof}

\begin{corollary}[An exact degree spectrum]\label{cor:degreespectrum}
In the noninvertible branch of Theorem~\ref{thm:triangular}, $h=0$ gives
coordinate degrees $(7,6,4)$. If $h\ne0$ has degree $m$, the coordinate
degrees are exactly
\begin{equation}\label{eq:extendeddegrees}
(2m+8,\ 2m+7,\ 2m+5).
\end{equation}
Consequently the possible maximum ordinary degrees in this noninvertible
class are precisely $7$ and every even integer at least $8$.
\end{corollary}
\begin{proof}
The coefficients of $z$ in the three coordinates of $F_{a,b}$ have
ordinary degrees $6,5,3$, respectively. Their leading terms are nonzero
for $a,b\ne0$. The substitution $z\mapsto z+y^2h(xy)$ adds terms of
degrees $6+(2+2m)$, $5+(2+2m)$, and $3+(2+2m)$. These are all strictly
larger than the original coordinate degrees, so they cannot cancel against
an original term. Their leading coefficients are nonzero, proving
\eqref{eq:extendeddegrees}. Every $m\ge0$ occurs by taking $h=t^m$.
\end{proof}

The first invariant coefficient is the important distinction here. The
coefficient of $v$ in $r$ is still required to be a scalar, not a
nonconstant polynomial in $t$. No assertion is made for
$r=R(t)+G(t)v$ with nonconstant $G$.

\section{Reduced parameter geometry and invertible degenerations}
\label{sec:geometry}

The formulas \eqref{eq:family} contain no negative powers of $a$ or $b$.
They therefore define a polynomial family on the entire parameter plane,
not just on its torus. The diagonal formula \eqref{eq:diagonal} should not
be used on the axes, where its individual scaling matrices are undefined.
The polynomial formulas are the correct extension there.

\begin{theorem}[The parameter plane and its boundary]\label{thm:plane}
In the finite coefficient space of normalized weighted maps of ordinary
degree at most seven with $p,q$ affine in $v$ and $r$ affine in $(t,v)$,
the following statements hold over $K$.
\begin{enumerate}[label=\textup{(\roman*)}]
\item The morphism $(a,b)\mapsto F_{a,b}$ is a closed embedding of
$\A^2_K$. Every member has Jacobian determinant $-1$.
\item Its torus $ab\ne0$ is exactly the normalized nonconstant-$r$ branch.
Its reduced Zariski closure is the entire embedded affine plane.
\item The invertible members of this plane are exactly its two coordinate
axes $V(ab)$. Explicitly,
\begin{align}
F_{0,b}(x,y,z)&=(z,\ y+9bxz,\ x),\label{eq:axisA}\\
F_{a,0}(x,y,z)&=\left(z-\frac{8a}{9}y^2,\ y,\ x\right).
\label{eq:axisB}
\end{align}
\item Inside this plane, the scheme-theoretic condition $r=1$ cuts out the
reduced nodal union $\Spec K[a,b]/(ab)$.
\end{enumerate}
\end{theorem}
\begin{proof}
The coefficients \eqref{eq:coeffrecover} recover $a,b$ by linear functions
on the ambient coefficient space. Every other coefficient is a polynomial
in those two functions by \eqref{eq:family}. Thus the image is the graph of
polynomial functions, a closed embedding.

The determinant is $-1$ on the torus by the classification and its explicit
solution. It is a polynomial in $a,b,t,v$, so the identity holds on the
whole plane. Equivalently, it can be checked directly from
\eqref{eq:family}. The torus is dense, and the classification identifies all
its field-valued nonconstant-$r$ points. This proves the closure statement.

On the axes, the formulas reduce to \eqref{eq:axisA} and \eqref{eq:axisB},
which are tame maps of the already classified constant branch. Off the
axes, \eqref{eq:newcollision} excludes invertibility. Finally, imposing
$r=1$ on the parameter ring imposes $(ab,ab^2)=(ab)$. The ideal $(ab)$ is
radical: it is the intersection $(a)\cap(b)$ in $K[a,b]$.
\end{proof}

\subsection{What degenerates, and what does not}

The Jacobian determinant stays $-1$ throughout this degeneration. The
ordinary degrees and the collision coordinates do not stay bounded in the
same manner. Fixing $b=1$ and letting $a$ approach zero over $\mathbb R$ or
$\mathbb C$, the coefficients tend to the automorphism $F_{0,1}$, whereas
the two source points in \eqref{eq:newcollision} have coordinates with
poles in $a$. In this explicit sense the displayed collisions escape to
infinity. No bounded limiting pair of distinct colliding points is asserted.

There are many other constant-$r$ maps, because $B\in t^2K[t]$ in
\eqref{eq:tamepq} is arbitrary. Only the two axes described above lie in the
closure of the nonconstant branch in this coefficient space. In particular,
passing from the nonconstant branch to its closure does not recover the
entire infinite-dimensional tame family.

\subsection{Rigidity over every reduced coefficient algebra}

Field-valued uniqueness can be strengthened without any degree cap on the
numerators.

\begin{corollary}[All-degree reduced rigidity]\label{cor:reduced}
Let $R$ be a reduced $\Q$-algebra. Suppose $p,q\in R[t,v]$ are affine in
$v$, satisfy the normalization and polynomiality conditions, and satisfy
$\J(p,q,1+t+v)=-1$. Then $p=p_*$ and $q=q_*$.
\end{corollary}
\begin{proof}
For every prime ideal $\mathfrak p$ of $R$, pass to the fraction field of
$R/\mathfrak p$. It is a characteristic-zero field, so
Theorem~\ref{thm:classification} forces $a=b=1$. Every coefficient of
$p-p_*$ and $q-q_*$ therefore belongs to every prime ideal of $R$, hence to
its nilradical. That nilradical is zero because $R$ is reduced.
\end{proof}

Thus fixing $r=1+t+v$ eliminates nonconstant families over reduced bases,
even when large ordinary degrees are allowed. It does \emph{not} eliminate
families over nonreduced bases. We now determine one such coefficient
scheme exactly.

\section{The normalized coefficient scheme is a double point}
\label{sec:scheme}

The preceding proof used degree arguments over fields. Nilpotent leading
coefficients need not satisfy those arguments. Scheme structure must
therefore be analyzed rather than inferred from the unique field-valued
solution.

\subsection{The precise finite coefficient problem}

Fix $r=1+t+v$ and ordinary degree at most seven. The most general normalized
numerators in our class are
\begin{align}
p={}&v+p_{20}t^2+p_{11}tv+p_{30}t^3+p_{21}t^2v
       +p_{40}t^4+p_{31}t^3v,\label{eq:scheme-p}\\
q={}&t+q_{01}v+q_{20}t^2+q_{11}tv+q_{30}t^3
       +q_{21}t^2v+q_{40}t^4.\label{eq:scheme-q}
\end{align}
The degree formulas in the proof of Corollary~\ref{cor:support} show that
this list is exhaustive. In particular, the $q_{40}$ term is allowed and
must not be set to zero beforehand.

Let the coefficient vector, in the following fixed order, be
\begin{equation}\label{eq:coordorder}
c=(p_{20},p_{11},p_{30},p_{21},p_{40},p_{31},
       q_{01},q_{20},q_{11},q_{30},q_{21},q_{40}).
\end{equation}
Let $I\subset\Q[c]$ be generated by all coefficients in $t,v$ of
$\J(p,q,1+t+v)+1$. There are eighteen nonzero coefficient polynomials.
Their complete expressions, ordering, and derivative matrix are supplied
in \texttt{verification.json}; the article fixes enough data below to
reconstruct and check the certificates independently.

\begin{theorem}[Exact double-point presentation]\label{thm:double}
The full coordinate algebra of this coefficient problem is
\begin{equation}\label{eq:dualnumbers}
\Q[c]/I\ \simeq\ \Q[\eps]/(\eps^2).
\end{equation}
Under this isomorphism $\eps=q_{21}-4$ and
\begin{equation}\label{eq:pointplus}
c=\cstar+u\eps,
\end{equation}
where
\begin{equation}\label{eq:cstar}
\cstar=\left(-\frac89,-2,\frac{28}{27},\frac43,-\frac8{27},-\frac8{27},
                     9,-8,-12,4,4,0\right)
\end{equation}
and
\begin{equation}\label{eq:u}
u=\frac1{144}(-7,-18,18,24,-8,-8,162,-198,-324,144,144,0).
\end{equation}
In particular, the entire ideal, not merely its radical, is generated by
$\eps^2$ and the eleven linear relations obtained from
\eqref{eq:pointplus} by omitting the tautological $q_{21}$ relation.
\end{theorem}

The proof has three components: uniqueness of geometric points, a rank
certificate for the tangent space, and a certificate excluding a
second-order lift. None requires a Gr\"obner-basis computation.

\subsection{Tangent space and its exact rank certificate}

Write the coefficient equations as $f_1,\ldots,f_{18}$ in the order of the
monomials
\begin{equation}\label{eq:monomialorder}
\begin{gathered}
t^7,\ t^6v,\ t^6,\ t^5v,\ t^5,\ t^4v,\ t^4,\ t^3v^2,\ t^3v,\ t^3,\\
t^2v^2,\ t^2v,\ t^2,\ tv^2,\ tv,\ t,\ v^2,\ v.
\end{gathered}
\end{equation}
Let
\begin{equation}\label{eq:M}
M=\left(\frac{\partial f_i}{\partial c_j}(\cstar)\right)_{i,j}.
\end{equation}
Direct rational calculation gives $Mu=0$. Delete the column corresponding
to $q_{21}$ and select the rows corresponding to
\begin{equation}\label{eq:minorrowmonomials}
t^7,\ t^6,\ t^5,\ t^4v,\ t^4,\ t^3,\ t^2v,\ t^2,\ tv^2,\ t,\ v.
\end{equation}
The determinant of that $11\times11$ submatrix is
\begin{equation}\label{eq:minorvalue}
-18\,874\,368\ne0.
\end{equation}
Thus $\operatorname{rank}M=11$ and $\ker M=\Q u$. This proves that the
scheme has a one-dimensional Zariski tangent space at its sole geometric
point. It does not yet determine the length of the point.

Because $f_i(\cstar)=0$ and $Mu=0$, the assignment
$c=\cstar+u\eps$ solves the equations over $\Q[\eps]/(\eps^2)$.
Since $u_{q_{21}}=1$, the resulting map from the coefficient algebra onto
the dual numbers is surjective. In particular, the tangent direction is
not merely a redundant parametrization of a reduced point.

\subsection{A nonzero obstruction at second order}

Define the quadratic residual vector
\begin{equation}\label{eq:H}
H=\bigl([\eps^2]f_i(\cstar+u\eps)\bigr)_{i=1}^{18}.
\end{equation}
Consider a possible lift to $\Q[\eps]/(\eps^3)$ of the form
$c=\cstar+u\eps+w\eps^2$. Its second-order equation is
\begin{equation}\label{eq:lift-obstruction}
Mw+H=0.
\end{equation}
There is an explicit row functional
\begin{equation}\label{eq:ell}
L=(-205,0,-46,0,-40,24,-8,0,4,0,0,0,0,0,0,0,0,0)
\end{equation}
satisfying
\begin{equation}\label{eq:obstructionpairing}
LM=0,\qquad LH=-\frac1{72}.
\end{equation}
Applying $L$ to \eqref{eq:lift-obstruction} produces a contradiction.
For a nonzero scalar multiple $\eta u$, the obstruction pairing is
$-\eta^2/72$, which is again nonzero over every characteristic-zero field.
There is therefore no second-order lift of any nonzero first-order tangent
direction. The matrix, vector, and pairing in
\eqref{eq:M}--\eqref{eq:obstructionpairing} are checked exactly by the
included script.

\subsection{From tangent obstruction to the entire scheme}

We record the small algebraic fact that turns these certificates into a
scheme-theoretic theorem.

\begin{lemma}[An isolated point with one obstructed tangent]
\label{lem:doublecriterion}
Let $k$ be algebraically closed and let $A$ be a finitely generated
$k$-algebra with exactly one geometric point and residue field $k$.
Suppose its tangent space has dimension one and no nonzero tangent
homomorphism $A\to k[\eps]/(\eps^2)$ extends to
$A\to k[\eps]/(\eps^3)$. Then $A\simeq k[z]/(z^2)$.
\end{lemma}
\begin{proof}
The Nullstellensatz implies that the radical of zero is the unique
maximal ideal $\mathfrak m$ \cite{nullstellensatz}. The algebra is
Noetherian. A finitely generated ideal consisting of nilpotent elements
is nilpotent: if its generators have nilpotence exponents $n_i$, every
product of more than $\sum_i(n_i-1)$ generators is zero. Thus some
power of $\mathfrak m$ is zero, and $A$ is a finite-dimensional local
$k$-algebra.

The cotangent space $\mathfrak m/\mathfrak m^2$ has dimension one.
Choose $z\in\mathfrak m$ whose residue spans it. Nakayama's lemma gives
$\mathfrak m=(z)$ \cite{nakayama}. Repeatedly writing an element as a
scalar plus $z$ times another element, and using nilpotence, shows that
$A$ is generated by $z$ over $k$. Since $k[Z]$ is a principal ideal
domain and the only root of the defining ideal is zero, it follows that
$A\simeq k[Z]/(Z^N)$ for some $N$. The nonzero cotangent space forces
$N\ge2$. If $N\ge3$, sending $Z$ to $\eps$ defines a homomorphism to
$k[\eps]/(\eps^3)$ with nonzero first-order part, contrary to the
hypothesis. Hence $N=2$.
\end{proof}

\begin{proof}[Proof of Theorem~\ref{thm:double}]
Work first over $\overline\Q$. The classification with $\beta=\gamma=1$
forces $a=b=1$, so the coefficient scheme has exactly one geometric point,
$\cstar$. Equations~\eqref{eq:minorvalue} and $Mu=0$ give tangent
dimension one, and \eqref{eq:obstructionpairing} excludes every nonzero
second-order lift. Lemma~\ref{lem:doublecriterion} gives a two-dimensional
coordinate algebra over $\overline\Q$.

Already over $\Q$, substitution $c=\cstar+u\eps$ defines a surjection
\[
\Q[c]/I\ \longrightarrow\ \Q[\eps]/(\eps^2).
\]
After extension to $\overline\Q$, this is a surjection between
vector spaces of dimension two and is an isomorphism. Field extension is
faithful, so its kernel over $\Q$ was zero. This proves
\eqref{eq:dualnumbers} with the specified coordinates.

The kernel of the explicit substitution map from the polynomial ring
$\Q[c]$ is the ideal generated by $\eps^2$ and the eleven affine-linear
relations $c_j-(\cstar)_j-u_j\eps$. The isomorphism identifies this kernel
with $I$, proving the exact ideal assertion.
\end{proof}

\subsection{The universal meaning of the answer}

\begin{corollary}[Solutions over arbitrary coefficient algebras]
\label{cor:functor}
For every $\Q$-algebra $R$, the solutions of
\eqref{eq:scheme-p}--\eqref{eq:scheme-q} with $\J=-1$ are in bijection with
\[
\{e\in R:e^2=0\}.
\]
The corresponding solution is $c=\cstar+ue$, and $e=q_{21}-4$.
\end{corollary}
\begin{proof}
Apply $\operatorname{Hom}_{\Q\text{-alg}}(-,R)$ to
\eqref{eq:dualnumbers}.
\end{proof}

In particular, every solution over a reduced ring is the constant one, but
a coefficient ring containing a nonzero square-zero element supports a
nontrivial infinitesimal deformation. The parameter is obstructed rather
than the initial segment of a moving curve of normalized field-valued
solutions. The square seen in the repository's finite coefficient
certificate is therefore genuine scheme structure, not just an artifact
of one Gr\"obner-basis presentation.

\subsection{The variable-parameter open scheme}

The nonconstant coefficient parameters can also be restored without
losing control of nilpotents. Let $\mathcal X^{\times}$ denote the
coefficient scheme obtained by using \eqref{eq:scheme-p}--\eqref{eq:scheme-q},
putting $r=1+\beta t+\gamma v$, and inverting $\beta\gamma$.

\begin{theorem}[Torus times the double point]\label{thm:torusdouble}
There is an explicit isomorphism
\begin{equation}\label{eq:torusdouble}
\mathcal X^{\times}\simeq
\Spec\Q[\beta^{\pm1},\gamma^{\pm1},\eps]/(\eps^2).
\end{equation}
It is an isomorphism of full schemes, not just of reduced varieties.
The square-zero coordinate can be taken to be
$\eps=q_{21}/(\beta\gamma)-4$.
\end{theorem}
\begin{proof}
Over any $\Q$-algebra where $\beta,\gamma$ are units, define
\begin{equation}\label{eq:scheme-normalization}
\bar p(t,v)=\gamma p(t/\beta,v/\gamma),\qquad
\bar q(t,v)=\beta q(t/\beta,v/\gamma).
\end{equation}
The third numerator becomes $1+t+v$. The two normalized origin
coefficients remain one. These substitutions preserve the allowed
monomial supports and ordinary degree bound, and have polynomial inverse
formulas over the Laurent coefficient ring. They preserve the determinant
condition: the row scalings contribute $\beta\gamma$ and the derivative
column scalings contribute $(\beta\gamma)^{-1}$.

Thus \eqref{eq:scheme-normalization} identifies the coefficient functor
with a choice of two units and a solution of the fixed normalized scheme.
Apply Theorem~\ref{thm:double} to that fixed scheme. The coefficient of
$t^2v$ in $\bar q$ is $q_{21}/(\beta\gamma)$, giving the stated
square-zero coordinate. All formulas are valid with nilpotent
coefficients, so this proves the scheme isomorphism.
\end{proof}

At every geometric point of this open scheme the reduced space has
dimension two and the tangent space has dimension three: two tangent
directions come from the torus and one from its square-zero thickening.
The remaining boundary problem is how this structure meets the axes,
where the Laurent normalization is no longer available.

\begin{remark}[The degree boundary matters]
Corollary~\ref{cor:reduced} is an all-degree statement. In contrast,
Theorem~\ref{thm:double} fixes the ordinary degree-seven coefficient space.
We have not proved that allowing higher-degree coefficients over
nonreduced rings introduces no extra tangent directions or higher-order
nilpotents. The field-valued classification alone does not imply that
stronger assertion.
\end{remark}

\section{Decision procedure, reproducibility, and formalization}
\label{sec:verification}

\subsection{A classification-based recognition procedure}

The theorem gives more than a determinant identity. For rational input
polynomials in the structural class, it yields a direct recognition
procedure for the Keller condition, together with an inverse or collision
certificate.

First check the prescribed weights monomial by monomial and extract the
unique invariant numerators as in Section~\ref{sec:coordinates}. Check
$\deg_vp,\deg_vq\le1$ and that $r$ is affine. Inputs outside these
restrictions are \emph{outside the procedure's scope}; they are not
rejected as non-Keller maps.

Within this scope, compute $p_v(0),q_t(0),r(0)$. If one is zero, the map
is not Keller by \eqref{eq:originJ}. Otherwise normalize the target
coordinates. Write $r=1+\beta t+\gamma v$. If $\beta=\gamma=0$, take
$B=p-v$ and $\lambda=\coeff{v}{q}$, and test the exact identities
\eqref{eq:tamepq} with $B\in t^2K[t]$. If exactly one of $\beta,\gamma$
is nonzero, the map is not Keller. If both are nonzero, recover $a,b$ from
\eqref{eq:abrecover} and compare with the fixed formula
\eqref{eq:family}. The classification proves both soundness and completeness
of these comparisons inside the stated scope.

For a positive result, return either the inverse
\eqref{eq:tameinverse} or the pair of source points in
\eqref{eq:newcollision}. To undo target normalization, compose an inverse
with the inverse target scaling; collision source points remain collision
source points, with their common target rescaled. This avoids a search over
all possible inverses or collision pairs.

The larger class of Theorem~\ref{thm:triangular} is recognized by first
computing its unique source shear. For $\gamma\ne0$, compute $h$ from
\eqref{eq:triangular-normal}, compose the source with $T_{-h}$, and apply
the affine-$r$ procedure. For $\gamma=0$, a nonconstant $R(t)$ is excluded
by that theorem. Inputs with nonconstant $v$-coefficient in $r$ remain
outside the classified scope.

For canonical sparse input, reading and comparing the coefficient lists
requires a number of field operations linear in the input support size,
apart from the cost of canonicalizing the input. This is an arithmetic
operation statement, not a claim of linear bit complexity for arbitrarily
large rational coefficients. An implementation should retain that distinction.

\subsection{What the companion script checks}

The file \texttt{verify\_results.py} works in exact rational and symbolic
arithmetic. It reconstructs the equations rather than trusting a stored
matrix. Its thirteen groups of checks cover:
\begin{enumerate}[label=\arabic*.]
\item The two-parameter determinant in both invariant coordinates and the
full three-variable Jacobian; the original-map normalization and collision;
diagonal equivalence, degrees, support, and the universal collision family.
\item The two degenerate axes, the two-sided inverse identity with a formal
arbitrary function $H$, the source-shear identities and sample degree
profiles, and the one-variable coefficient identities and explicit endpoint
of the classification proof.
\item The invertible variable-parameter scheme normalization, the eighteen
coefficient equations, the rank minor, the tangent
vector, the cokernel functional, the second-order obstruction, and the
vanishing of every coefficient equation under the dual-number presentation.
\end{enumerate}
The last check proves one ideal containment by direct substitution. The
reverse containment is justified in the article by the classification,
rank, and obstruction argument; the script does not silently label a
one-sided ideal check as equality.

The run included with this report used SymPy~1.14.0 and passed all thirteen
check groups. The JSON certificate includes the full rational matrix and
both row and column orderings. No numerical tolerance, random test, or
floating-point approximation is used. These checks are reproducibility
evidence for the finite algebra; they do not replace the universal
one-variable degree proofs.

\subsection{A realistic proof-assistant boundary}

The new theorems in this report have not been compiled in Lean or Rocq.
The existing repository formalizations of the original map are separate
from the new classification. A sensible formalization order is the
following dependency chain:
\[
\begin{gathered}
\text{weighted monomial encoding and localization identity}\\
\Downarrow\\
\text{constant weighted Wronskian lemma and cubic--square lemma}\\
\Downarrow\\
\text{three coefficient equations and exhaustive case split}\\
\Downarrow\\
\text{explicit classification, inverse, and collision theorems}\\
\Downarrow\\
\text{finite matrix certificates and double-point criterion}.
\end{gathered}
\]
The main classification needs polynomial degrees, characteristic-zero
derivatives, a rational-function constant lemma, and unique factorization
in one variable. It does not require the heavy geometric machinery used
in the literature to study unrestricted Keller maps. The double-point
statement additionally needs finite-type algebra, the Nullstellensatz,
Nakayama's lemma, and faithful scalar extension, or an alternative direct
ideal-membership certificate.

One attractive alternative for formalizing the last step is to produce
explicit polynomial combinations expressing the twelve simple generators
of the double-point ideal in terms of the eighteen original equations.
The present argument proves these combinations exist but does not supply
their expanded coefficients. Such certificates could isolate the formal
scheme result from a large commutative-algebra library.

\section{Further research questions and proposed next steps}
\label{sec:future}

The following questions are proposed by the present results. Unless
explicitly attributed, they are research directions formulated here, not
claims that a published source listed them as open. They are ordered to
separate immediately testable extensions from much broader classification
problems.

\subsection{Allow quadratic dependence on the second invariant}

Keep $r$ affine but allow $\deg_vp,\deg_vq\le2$. Does the nonconstant
branch admit genuinely new diagonal-equivalence classes? After the change
$s=r$, the determinant becomes a polynomial of higher degree in $s$.
The highest coefficient should impose a weighted differential relation on
the leading coefficients, but the lower equations can couple several
common factors. The first concrete task is to derive the complete
coefficient system and identify whether an analogue of
Lemma~\ref{lem:wronskian} gives an a priori degree bound. A finite search
would be useful evidence, but an all-degree conclusion needs such a bound
or another structural argument.

\subsection{\texorpdfstring{Allow a nonconstant second-invariant coefficient in $r$}{Allow a nonconstant second-invariant coefficient in r}}

Theorem~\ref{thm:triangular} settles arbitrary polynomial $t$-dependence
when the coefficient of $v$ in $r$ is constant. The next case is
$r=R(t)+G(t)v$ with nonconstant $G(t)$. Must $G$ be constant under suitable
Keller and polynomiality hypotheses, or do genuinely new normal forms
appear? The substitution $s=r$ now requires division by $G$ and is not a
global polynomial coordinate change. Track the behavior at the zeros of
$G$ and determine whether the constant-Wronskian argument has a localized
analogue. Allowing a quadratic term in $v$ is a further, distinct extension.

\subsection{Does the nonreduced structure stabilize with degree?}

For each $d\ge7$, form the normalized coefficient scheme with
$r=1+t+v$, $p,q$ affine in $v$, and lifted ordinary degree at most $d$.
Its reduced scheme is a point by Corollary~\ref{cor:reduced}. Is the full
scheme always the same double point? Or can high-degree coefficients
support new nilpotent directions invisible at degree seven? A first step
is to linearize the three identities in $(A,B,C,D)$ around the explicit
solution without truncating their degrees, solve the resulting polynomial
linear differential system, and then examine its obstruction maps.
This is a particularly well-targeted extension of the present work.

\subsection{Scheme-theoretic gluing at the invertible boundary}

Theorem~\ref{thm:torusdouble} settles the full degree-seven coefficient
scheme on $\beta\gamma\ne0$: it is the parameter torus times a double
point. Determine how this thickening meets the constant-$r$ tame locus
where that normalization ceases to be invertible. Does its scheme-theoretic
closure acquire embedded components or higher nilpotent order on either
axis or at their intersection? The reduced affine-plane closure and the
nodal boundary in Theorem~\ref{thm:plane} do not decide this gluing problem.
A concrete first step is to saturate the coefficient ideal by
$\beta\gamma$, then study its specialization on each axis with explicit
ideal-membership certificates.

\subsection{Other weights and other Wronskian exponents}

Our determinant produces a cubic--square relation because the first two
weights are $2$ and $1$. For other mixed-sign weight triples, analogous
highest-coefficient equations may have the form
$mCA'-nAC'=0$. Investigate the resulting common-power factorizations and
constant weighted Wronskians. The interesting issue is not just repeating
the calculation with new constants: the polynomiality inequalities at the
coordinate hyperplane may either permit or rule out the required
factorizations. A systematic analysis could delineate entire rigid classes
of equivariant maps.

\subsection{Compare lower-degree maps outside this symmetry class}

The current literature already includes a three-dimensional degree-four
example \cite{gao}; the existence of such an example is therefore not
posed as a new open problem here. Instead, extract its invariant or
geometric structure and determine precisely which hypothesis of our
classification fails. Compare support size, rational coefficient height,
and the cost of a formal collision certificate. This would help replace a
single ``simplest map'' question by explicit, noninterchangeable complexity
criteria.

\subsection{Arithmetic optimization within the diagonal orbit}

Over $\Q$, the classification makes every nonconstant map in this class
explicitly diagonal-equivalent. One can now optimize primitive coefficient
height, denominator size, or integral collision height within that orbit.
The parameters $a,b$ and the additional unnormalizing target scalings make
these concrete arithmetic optimization problems. Prove lower bounds for a
specified normalization rather than relying on finite samples of small
rational parameters.

\subsection{Direct ideal certificates and kernel verification}

Generate explicit membership certificates for the twelve simple generators
in Appendix~\ref{app:presentation}, and check them with polynomial
normalization in Lean or Rocq. This would turn the deformation theorem into
a compact finite certificate after the all-degree classification has been
formalized. The target should be a transparent theorem boundary: no hidden
use of a computer-algebra claim that a Gr\"obner basis is complete.

\subsection{Positive characteristic as a controlled contrast}

The displayed family and finite identities have meaningful reductions at
many primes, but the all-degree proof does not automatically survive.
In positive characteristic, nonconstant $p$th powers have derivative zero,
and the weighted degree coefficients can vanish. Classify which parts of
the theorem persist under explicit degree bounds and which fail through
Frobenius phenomena. Separately determine the primes at which the tangent
rank or obstruction changes. Reducing a characteristic-zero certificate
is not sufficient to claim the same unrestricted classification.

\section{Conclusion}

The repository's finite sparsity problem admits a stronger structural
answer in its affine-invariant class. Constant third numerator gives all
and only the explicit tame maps \eqref{eq:tamepq}. Nonconstant third
numerator forces both invariant coefficients to be nonzero and forces the
entire map into the diagonal orbit \eqref{eq:family}, with degree seven and
sixteen terms. An explicit source shear extends the classification to
$r=R(t)+\gamma v$ and gives the exact degree spectrum
$\{7,8,10,12,\ldots\}$. The proof is an all-degree argument, not an
extrapolation from a coefficient search.

The parameter-space results explain two different notions of rigidity.
Before fixing the third numerator, the nonconstant branch has an affine
plane closure with invertible coordinate axes. After fixing $r=1+t+v$,
all reduced families are constant, yet the degree-seven coefficient scheme
retains exactly one square-zero parameter. Its tangent direction has an
explicit nonzero second-order obstruction. Restoring the two nonconstant
parameters gives the torus times this double point. These distinctions are essential
both for further searches and for a trustworthy formalization.

\appendix
\section{An explicit presentation of the coefficient ideal}
\label{app:presentation}

For readers who prefer generators to the vector formula
\eqref{eq:pointplus}, Theorem~\ref{thm:double} says that $I$ is exactly
the ideal generated by
\begin{equation}\label{eq:idealexplicit}
\begin{aligned}
&q_{40},\qquad (q_{21}-4)^2,\\
&144p_{20}+7q_{21}+100,\\
&8p_{11}+q_{21}+12,\\
&216p_{30}-27q_{21}-116,\\
&6p_{21}-q_{21}-4,\\
&54p_{40}+3q_{21}+4,\\
&54p_{31}+3q_{21}+4,\\
&8q_{01}-9q_{21}-36,\\
&8q_{20}+11q_{21}+20,\\
&4q_{11}+9q_{21}+12,\\
&q_{30}-q_{21}.
\end{aligned}
\end{equation}
These equations apply over every $\Q$-algebra, not just over fields.
Over a field the square forces $q_{21}=4$, and all remaining coefficients
are the point \eqref{eq:cstar}. Over the dual numbers, $q_{21}=4+\eps$
produces exactly the tangent vector \eqref{eq:u}.

As a useful independent check of the equation reconstruction, the
coefficients corresponding to the last three monomials $t,v^2,v$ in
\eqref{eq:monomialorder} are, respectively,
\begin{align*}
f_{16}&=-p_{11}+2p_{20}q_{01}-2q_{20}-2,\\
f_{17}&=2p_{11}q_{01}-3q_{11},\\
f_{18}&=p_{11}q_{01}+q_{01}-q_{11}-3.
\end{align*}
For the first two monomials they are
\[
f_1=-5p_{31}q_{40}-4p_{40}q_{40},\qquad
f_2=-9p_{31}q_{40}.
\]
The remaining equations are recovered by the same determinant expansion;
the complete list is included in the JSON certificate.

\section{Small exact certificates for the deformation calculation}
\label{app:certificates}

The primitive integral tangent vector is $n=144u$. To keep the second-order
residual integral, define
\[
\widehat H_i=[e^2]f_i(\cstar+en)=144^2 H_i.
\]
In the ordering \eqref{eq:monomialorder}, the nonzero data and the
cokernel functional are as follows. Zeros are included to make the ordering
unambiguous.
\begin{center}
\begin{tabular}{@{}rcrr@{}}
\toprule
Row & Monomial & $\widehat H_i$ & $L_i$\\
\midrule
1 & $t^7$ & 0 & $-205$\\
2 & $t^6v$ & 0 & 0\\
3 & $t^6$ & 0 & $-46$\\
4 & $t^5v$ & 0 & 0\\
5 & $t^5$ & 1008 & $-40$\\
6 & $t^4v$ & 1872 & 24\\
7 & $t^4$ & $-4140$ & $-8$\\
8 & $t^3v^2$ & 864 & 0\\
9 & $t^3v$ & $-9504$ & 4\\
10 & $t^3$ & 2772 & 0\\
11 & $t^2v^2$ & $-5184$ & 0\\
12 & $t^2v$ & 12636 & 0\\
13 & $t^2$ & 1620 & 0\\
14 & $tv^2$ & 9720 & 0\\
15 & $tv$ & $-2592$ & 0\\
16 & $t$ & $-2268$ & 0\\
17 & $v^2$ & $-5832$ & 0\\
18 & $v$ & $-2916$ & 0\\
\bottomrule
\end{tabular}
\end{center}
The obstruction is visible in four integer products:
\[
(-40)(1008)+24(1872)+(-8)(-4140)+4(-9504)=-288.
\]
Dividing by $144^2$ gives $LH=-1/72$.

For the nonzero minor in \eqref{eq:minorvalue}, the zero-based row indices
in the JSON matrix are
\[
(0,2,4,5,6,9,11,12,13,15,17),
\]
and the zero-based column indices are
\[
(0,1,2,3,4,5,6,7,8,9,11).
\]
Thus an independent checker only needs the reconstructed derivative
matrix, one $11\times11$ determinant, and the two matrix-vector products
$Mu$ and $LM$. The companion JSON stores the matrix itself as rational
strings for convenient cross-checking.

\section{Reproduction and source ledger}
\label{app:reproduction}

The supplied archive contains the following files:
\begin{center}
\begin{tabular}{@{}p{0.36\textwidth}p{0.56\textwidth}@{}}
\toprule
File & Purpose\\
\midrule
\texttt{article.tex} & Self-contained LaTeX source, including bibliography.\\
\texttt{article.pdf} & Compiled article.\\
\texttt{verify\_results.py} & Reconstructs and checks all finite identities.\\
\texttt{verification.json} & Exact equations, matrix, and certificate data.\\
\texttt{verification.txt} & Output of the successful verification run.\\
\texttt{README.md} & Build commands, scope, and trust boundary.\\
\texttt{requirements.txt} & Reproducible SymPy dependency.\\
\bottomrule
\end{tabular}
\end{center}
Build the article with two runs of \texttt{pdflatex article.tex}, or with
\texttt{latexmk -pdf article.tex}. Run the checks with
\begin{quote}\ttfamily
python verify\_results.py --output verification.json
\end{quote}
The verification does not download or execute any repository code. The
repository is a cited mathematical source; its original map and the
relevant formulas are reconstructed in the script. The source ledger is:
\begin{center}
\begin{tabular}{@{}p{0.30\textwidth}p{0.61\textwidth}@{}}
\toprule
Source category & Role in the report\\
\midrule
Repository map & Original map, its known collision, and the finite-search
baseline; not new claims of this article.\\
External research & Context for graded maps and examples outside our class;
not dependencies of the classification proof.\\
Standard algebra & Nullstellensatz and Nakayama's lemma in the passage from
an isolated obstructed tangent to a double point.\\
This article & All-degree classification, family closure, reduced rigidity,
and the full degree-seven coefficient-algebra identification.\\
\bottomrule
\end{tabular}
\end{center}

\begin{thebibliography}{9}
\small
\bibitem{repo-main}
Vladimir Reshetnikov, \emph{ProveIt: Algebra/JacobianConjecture/README.md}.
Repository snapshot \texttt{e21766d04c2b8a9b2cdba0cd43e563024b3bd1b9},
inspected September 24, 2026.
\url{https://github.com/VladimirReshetnikov/ProveIt/blob/e21766d04c2b8a9b2cdba0cd43e563024b3bd1b9/Algebra/JacobianConjecture/README.md}.

\bibitem{repo-research}
Vladimir Reshetnikov, \emph{ProveIt: Algebra/JacobianConjecture/Research/README.md}.
Same repository snapshot and inspection date.
\url{https://github.com/VladimirReshetnikov/ProveIt/blob/e21766d04c2b8a9b2cdba0cd43e563024b3bd1b9/Algebra/JacobianConjecture/Research/README.md}.

\bibitem{repo-sparsity}
Vladimir Reshetnikov, \emph{verify\_equivariant\_sparsity.py}.
Same repository snapshot and inspection date.
\url{https://github.com/VladimirReshetnikov/ProveIt/blob/e21766d04c2b8a9b2cdba0cd43e563024b3bd1b9/Algebra/JacobianConjecture/Research/verify_equivariant_sparsity.py}.

\bibitem{shaska}
T.~Shaska, \emph{Graded Keller maps and the Jacobian Conjecture},
arXiv:2607.20210, 2026. Version 2 dated July 25, 2026;
consulted September 24, 2026.
\url{https://arxiv.org/abs/2607.20210}.

\bibitem{gao}
Shuhong Gao, \emph{Counterexamples to the Jacobian conjecture in dimensions
greater than two}, arXiv:2608.00222, 2026. Version 1 dated July 31, 2026;
consulted September 24, 2026.
\url{https://arxiv.org/abs/2608.00222}.

\bibitem{nullstellensatz}
The Stacks Project Authors, \emph{The Stacks Project},
Theorem 10.34.1, Hilbert Nullstellensatz, Tag 00FV.
Consulted September 24, 2026.
\url{https://stacks.math.columbia.edu/tag/00FV}.

\bibitem{nakayama}
The Stacks Project Authors, \emph{The Stacks Project},
Lemma 10.20.1, Nakayama's lemma, Tag 00DV.
Consulted September 24, 2026.
\url{https://stacks.math.columbia.edu/tag/00DV}.
\end{thebibliography}
\end{document}
